Modern datacenter applications like social networking, financial trading systems, and online video conferencing operate under stringent Service Level Objectives (SLOs).
These SLOs span from tail latency targets to deadlines, and are critical for ensuring user satisfaction and meeting business goals.

The diversity of SLOs necessitates tailored scheduling policies to meet specific application requirements.
% The diversity in SLOs and workload characteristics across applications necessitates the need for tailored scheduling policies to meet specific application requirements.

However, designing and implementing such policies for application-specific SLOs is a complex task, requiring deep expertise in both the application domain and low-level scheduling mechanisms.
% Monolithic kernels provide only a handful of scheduling policies, prioritizing generality over per-application requirements.
% Prior work~\cite{humphries2021ghost,linux-sched-ext} improves kernel scheduling flexibility through extensible scheduling classes, but remains constrained by the kernel's millisecond-scale scheduling granularity.
% To avoid kernel scheduling overhead, recent research has proposed highly specialized dataplane operating systems (OSes) for microsecond-scale workloads.
% Yet, despite their performance, these systems typically hardcode a single scheduling model for all workloads, making per application customization impractical.

In this paper, we present \sysname, a scheduling synthesis framework that automatically generates schedulers from application-defined policies and intent.

% Our key insight is that the rigidity of state-of-the-art dataplane OSes stems from \textit{conflating these three concerns}.
% Specifically, existing systems treat tail latency as the sole scheduling intent, realize it through a fixed policy, typically a hybrid of First Come First Serve (FCFS) and Processor Sharing (PS), and enforce that policy in the runtime scheduler.
% Such conflation leads to inflexible scheduling, as applications have no authority over the scheduling intent or policy, and the OS arbitrarily enforces a one-size-fits-all policy that may not align with the application's requirements.

% % \sysname relies on three key mechanisms to empower applications to regain control over scheduling.
% % First, \sysname provides \textbf{Policy Engine} as the primary interface for applications to define their scheduling policies.
% % % The Policy Engine exposes a set of  
% % Second, \sysname enforces application-defined scheduling policies through a trusted scheduler, \textbf{MesoCore}, which runs inside the \textbf{Shadow Ring}. 
% % Finally, \sysname introduces the \textbf{Interrupt Express} subsystem in the Kernel to securely multiplex low-level hardware between the MesoCore and the Kernel Scheduler.

% \sysname enables application-defined scheduling through three key mechanisms.
% First, the \textbf{Policy Engine} provides a declarative interface that decouples scheduling policy and intent from enforcement.
% Applications express policies via a \textit{rank} function and specify SLOs such as tail latency targets or deadlines, without concern for how these are realized at the hardware level.
% Second, the \textbf{Synthesis Engine} bridges this high-level specification with low-level execution.
% It statically analyzes the rank function to determine rescheduling conditions and maps them to concrete hardware events (e.g., I/O completions, timer expirations, queue state changes).
% It then generates a scheduling handler that recomputes priorities and selects the next task when these conditions are met and the hardware events occur.
% Finally, \sysname enforces the synthesized implementation through a trusted scheduler, \textbf{MesoCore}. 
% MesoCore runs inside the \textbf{Shadow Ring} and directly interacts with low-level mechanisms such as user interrupts and hardware timers to achieve microsecond-scale enforcement, while preserving isolation from untrusted applications.
% % Crucially, rather than exposing vulnerable hardware primitives like raw APIC timer or user-level Inter-Processor Interrupts (IPIs), the Policy Engine allows applications to safely express complex scheduling semantics through a secure, shared-memory dispatch descriptor.
% % Second, \sysname enforces application-defined scheduling policies through a trusted scheduler, \textbf{MesoCore}, which runs inside the \textbf{Shadow Ring}.
% % Rather than exposing vulnerable hardware primitives directly to user space, MesoCore securely parses the application's descriptor and translates these logical decisions into bare-metal actuations.
% % Specifically, MesoCore programs core-local APIC timers and issues IPIs to trigger strict preemption at applications' command \yihan{too vulgaris}.
% % Finally, \sysname introduces the \textbf{Interrupt Express} subsystem in the Kernel to securely multiplex low-level hardware between MesoCore and the Kernel Scheduler.